\chapter[Sermon 32. The Fifth Seal Is Opened, and the Persecution of the Faithful Set Before Our Eyes, and Also the State of Martyrs in Another World.]{The Fifth Seal Is Opened, and the Persecution of the Faithful Set Before Our Eyes, and Also the State of Martyrs in Another World.}
\label{sermon:032}
\markright{Sermon 32. The Fifth Seal Is Opened, and the Persecution of the ...}

And when he had opened the fifth seal, I saw under the altar the souls of those who had been killed for the word of God, and for the testimony they had borne. They cried out with a loud voice, saying, “How long, O Lord, holy and true, will you delay to judge and to avenge our blood on those who dwell on the earth?” And white robes were given to each of them, and it was said to them that they should rest for a little season, until the number of their fellow servants and brethren, who were to be killed as they had been, should be fulfilled.

The fifth seal being opened by the Lamb, he exhibits to our eyes, or rather compels us to see, the continual persecutions of the church; and he shows us diligently what is the state of those who die in persecutions. Truly, the Lord Christ sends forth ministers and preachers for the salvation of men. And they, unthankful, overwhelm the faithful messengers of God with all kinds of injuries, and at length most cruelly slay them. Concerning this matter, the talk of men among themselves is diverse; but the very Son of God at this present boldly instructs his church, declaring what the godly shall suffer.

\index{Antichrist}\index{Patmos}\index{Cyprian, Saint}\index{Eusebius of Caesarea}\index{Rome}And first in expounding the same, we shall speak generally of the persecutions, by which both the ministers and the entire faithful church are diversely afflicted. The Lord Christ has shown us before in the Gospel many things concerning the persecutions to come, truly that he might prepare the minds of all the faithful for battle and patience. The passages are in Matthew 10:1-24, Luke 12:1-21, and John 14:15-16. And also the Acts of the Apostles tell of many things which the godly suffered in that most holy primitive Church. Would anyone have been considered sane if, at that time, they had said: It appears that the apostolic church is not the church, for that it is subject to all the mockeries, injuries, and slaughters of all men? Why, then, do we not acknowledge at this day that those are greatly deceived who measure the church by the outward peace and tranquility of things? Paulus Orosius in the seventh book of histories recounts ten grievous persecutions raised against the church from the time of the Apostles until the emperor Constantine, a time that did not fully accomplish the space of three hundred years. The first was stirred up by Nero, a monstrous man, of whom also Tacitus mentions in his Chronicles. This man put out of the way Peter and Paul, the most holy Apostles of Christ. The second destruction of the church came with Domitian, who in that persecution most grievously afflicted both our John and the whole church also, and when he was brought to Rome, banished him to the Isle of Patmos. The third was raised by Trajan, of whom Pliny, governor of Asia, makes mention in the tenth book of Epistles. In this persecution Ignatius, a holy bishop, was cast to the beasts and devoured. And Marcus Aurelius molested the church with the fourth persecution, and consumed with fire Polycarp, a bishop most worthy. Septimus Severus moved the fifth persecution, which Eusebius pursues in the sixth book of the \textit{Ecclesiastical History}. Julius Maximinus killed Pamphilus, a martyr, and Sextus raged cruelly against the church. And Decius Trajan began the seventh persecution, and executed very many who professed Christ. And Licinius Valerian, emperor, beheaded Cyprian, the good bishop of Carthage, and was the eighth persecutor of the church. Aurelian began the ninth persecution, which he but little advanced, for God, most just, took him away immediately. But Diocletian and Maximian shed more Christian blood than any other of the Roman emperors. Read, I pray you, the beginning of the eighth book of the \textit{Ecclesiastical History} of Eusebius. Compare those things with our time, and judge and conjecture what will shortly come to pass, and what our state will be. Persecutions are again renewed after Constantine, under Constantius and Julian. But the most terrible and grievous of all have boiled up under Antichrist, and have endured now for the space of five hundred years and more. What is done at this day, all the world sees. The ground is wet with the blood of martyrs, which things John foresaw.

\index{Daniel (Book of)}And the causes of persecution arise partly from the rule of Christ, which opens here the fifteenth Seal, and partly from men. The Lord sends to his people the cross and suffering to quicken those who are slow, to cleanse those covered with rust, and to refine corrupt gold. For so the Scripture declares in the eleventh chapter of Daniel and the Apostle Peter in the fourth chapter. Christ therefore does not destroy, but tests, and permits many things to tyrants against the Church. Godly men also bring upon themselves the Lord’s heavy hand, while indeed they believe rightly in the Son of God and depend only on him, but nevertheless are entangled with various and evil affections and commit acts that do not befit them. This may be seen declared at length in the beginning of the eighth book of the \textit{Ecclesiastical History} of Eusebius, which I recently cited. And the tyrants who persecuted had a different motivation: as Sennacherib and Antiochus did, than our bishops and princes do today. For these now are moved by hatred of religion and are spurred forward by Satan. They will by all means have their idolatrous religion maintained and the religion of the Gospel utterly destroyed. They cannot tolerate having their idols or other sins reproved; for this cause they are enraged at the faithful and those who frankly speak against and blame their idols and wickedness. And thus does the persecution arise, boil up, and proceed.

When the faithful see such increase [?] and feel deeply oppressed, they wonder how long the Lord will overlook this. Many cry out that the Lord is neglecting his affairs. The Lord seems to deal unjustly with his servants; he appears utterly to forget them. There is no doubt that many, by murmuring, offend the Lord grievously. Now therefore, we are taught to have hope and patience.

And at this present, heaven is opened to us, and showed us what is to be seen, namely, the souls of those who have been slain in persecutions, and their state is declared. Moreover, that God does not forget to avenge; why also he delays it, and for how long. These things are spoken for the consolation of all the faithful who are now afflicted by persecution. Far different things are exhibited here than those painters, instructed or rather corrupted by monks and friars, have shown us: namely, a great company of monks and nuns covered in heaven with our Lady’s cloak, as though the greatest part of them were to be saved. John shows us never a friar, but rather many martyrs, whom the friars at this day make before other men. Hereof, therefore, as of the doctrine of truth, we shall learn what state or degree is most plentiful in heaven, not that we should think no man but only martyrs are to be saved (for as many as truly believe in Christ, and crucify their flesh with its desires, shall be associated with the holy martyrs, and rejoice with Christ forever), but that chiefly the holy martyrs are saved, whom the mad world supposes to be lost.

But all things here must be examined by us with the greatest diligence. For this passage is both very clear and full of most wholesome doctrines. First, John sees and shows us, as it were pointing with his finger, the souls, and those of those who were slain—that is, the spiritual and immortal substances which remain alive when the body is lost and consumed. The body may be killed, but the soul cannot be killed. Our Savior has vividly expressed this in Matthew 10. In Luke 12, he says, “Do not fear those who kill the body, and after that can do nothing more,” and so on. Therefore, tyrants could rightly kill the bodies of martyrs; they had no power over their souls. This passage clearly demonstrates that the souls of men are not only immortal but also living or watchful, not sleeping, and remain and live in heaven. For some believe that the souls, when they depart from the body, sleep—which is a most vain notion.

Now the reason for the martyrs’ deaths is also shown: they were put to death for the word of God and for the testimony they had. They were not put to death for their wickedness or evil deeds, but for the true religion, by which they confessed and preached that word of God which was in the beginning and was made flesh, and the Gospel which they had committed to them, the testimony of God and eternal life, which also they ministered and preached. Of the word of God and the testimony of Jesus Christ I have spoken in the first chapter. For no other cause are countless bishops, kings, and princes slain today. If they were adulterers, usurers, blasphemers, and wicked doers, they would be regarded with some measure of tolerance; but now, because they profess the only Son of God and preach the Gospel, they are murdered without mercy. Here we have certainly defined who are true martyrs in deed, not those who suffer torments, but those who are tormented for God’s word. For the cause makes the martyr.

\index{Aquinas, Thomas}But where are the souls of those who are slain for the word of God? They are shown to us under the Altar; the Altar, as in the eighth chapter, is set in heaven, before the throne of God. Therefore, the souls of all the saints are in heaven, before the throne of God, which was also signified before in the vision of the twenty-four Elders. The Lord has said also, “Where I am, there shall also my servant be.” But the Lord is in heaven; therefore, the souls of the faithful, whose bodies have been slain, or buried without slaughter, are nowhere else but in heaven. Nevertheless, it does not lack a singular mystery that they are laid under the Altar, as under a shadow, through whose benefit the souls may be well at ease. I told you before, and here again repeat, that the altar signifies Christ. For he is also the golden altar, intercessor, and propitiation for our sins. Through the propitiation and mediation of Christ, we are received into celestial joys. And Christ is our life and salvation. Under him we lie hid, as under a cover or a shadow. Thomas Aquinas, expounding this passage from John, says that by the altar is signified Christ, in whom and by whom we should offer to the Father whatever good we do; and through him, whatever is pleasing to God is made acceptable. Under this Altar, namely under Christ, are the souls, not only in the state of life (that is, while we live here on earth) but also in the state of our homeland (that is, in heaven), as under him of whom they are covered, as under a shadow against all evil. Thus says Thomas. But I suppose that there is another thing also signified, that martyrs are made conformable to the altar, that is, to the passion of Christ, and therefore to rest now under the Altar, Christ. For those who are partakers with him in passion, do communicate also with him in glory. For just as the bosom of Abraham is called a receptacle, and that port and haven of salvation, into which the souls of those who had the faith of Abraham are received, so we understand the altar to be a place of blessedness in heaven, wherein they rest who with true faith have acknowledged Christ as the altar, propitiation, sanctification, and satisfaction; and have moreover, in suffering, offered themselves to God in Christ, through patience, an acceptable sacrifice to God. Under this Altar was gathered the first martyr, Abel; and after him, as many as have died for the faith, and shall be gathered, whoever in bearing the cross, through tribulation, enter with Christ into glory.

Now it is also declared what they do under the altar. I say, the very martyrs cry out, not the beasts, as they have done until now, and they cry out with a loud voice. Let no one imagine that the blessed souls in heaven are complaining, are sorrowful, accuse, or are troubled. These things are feigned for another purpose, so that we may understand that God does not forget his people, that he does not withhold vengeance, that he sees, feels, and regards the injuries and deaths of his servants. For when vengeance does not follow immediately, many think that God is asleep and has no regard for them. Therefore, we hear that the holy martyrs cry out, and with a loud voice. He appears to have alluded to that same in Genesis 4: “The voice of your brother’s blood cries to me, for vengeance.”

\index{Hebrews, Epistle to the}For the theologians call certain sins “crying,” as those which are red in the Scriptures cry out to God, as is presently the shedding of blood; the sin of Sodom in Genesis 9; the oppression of widows and orphans in Exodus 22; wages for work detained, Deuteronomy 24; and James 5. How long so ever therefore God delays vengeance, even if it is for many years, the blood of the just is not forgotten before God. Paul in Hebrews 12 cries out and says that the blood of Abel speaks. In Luke 18, the Lord says that the afflicted cry out both day and night for deliverance. Would God they would weigh these things, whose feet are swift to shed blood. God did not show mercy to his people in times past because much innocent blood was shed among them through the means of Manasseh, their king, as appears in the 4th book of Kings. Therefore, dear brethren, let us consider well at this day what we do, and let us not shed innocent blood rashly.

\index{John, Saint}\index{Augustine, Saint}Certainly the words are expressed of Saint John, which the martyrs cried to the Lord: how long, say they, Lord, which art holy and true, \&c. They put God in remembrance, not as ignorant or inconstant, but as knowing and steadfastly mindful of holiness and truth. For in as much as the Lord is holy, he hates all profane and unclean persons, and spares them not. For as much as he is true, he maintains and defends his chosen, and punishes and oppresses his enemies as he has promised by his word. Sins therefore you are such, say they, O God, why doest thou not judge and avenge our blood, of them which in earth, as in their kingdom, exercise tyranny and oppress every good man? All this signifies no other thing than that God for his own sake, which is holy and true, will never forget the injuries of his servants. Therefore we understand these things to be spoken by a figure called Prosopopeia: that is, the feigning of a person; not that the saints in Heaven do expostulate with God, but that we by such a figure might understand that God has care of martyrs, because he is holy and true. Saint Augustine in the 68th question upon the New Testament: Saying the Lord, says he, has taught us to pray for our enemies, what is the cause that the souls of those that are slain cry out as doth the blood of Abel, and require that they may be avenged? And he makes answer: Saints be not impatient, that they should urge that thing to be done now, which they know shall come to pass in the time prefixed, which neither can be prevented, nor yet delayed; but by this saying he would show, how God will avenge the blood of his servants, least because he seems now so patient, that wicked war should be thought unpunished, which is made against the saints: that both he might drive fear into them that persecute the servants of God, and might also exhort the sufferers unto patience. Thus says he. And this in deed seems the plainest sense of all others, especially if we consider the things that follow in the Lord’s answer, and it was said unto them that they should rest, \&c.

\index{Arethas of Caesarea}\index{Primasius}\index{Resurrection}Primasius, Bishop of Utica, explaining this passage in the writings of Saint John, states that it should not be understood as the saints seeking vengeance with a carnal understanding and stubbornness, for we know that through the abundance of charity even enemies are beloved in this case. Rather, it is evident that they prayed against the kingdom of sin, and earnestly desired the things of that kingdom, as we say, “thy kingdom come.” It is not lawful to think that they would desire anything against the pleasure of God, since their desires depend upon his will, and so forth. And Saint Gregory asks, what is it that the souls request when they ask for vengeance, but that they desire the last day of judgment and the resurrection of slain bodies? Arethas notes here also from the commentaries of Saint Andrew, Bishop of Caesarea: moreover, the saints appear through this to wish for the end of the world. Therefore, they are commanded to patiently abide until the fulfillment of their brethren, lest they should be fulfilled without them, as the holy Apostle says.

\index{Apocalypse (Book of)}However, Thomas Aquinas in his explanation of the Apocalypse shows that vengeance is required of God in two ways. First, with an evil and malicious affection, which Scripture utterly condemns. Secondly, by a zeal for righteousness, and according to God’s will, judgment is required against those who are incurable. After this, he adds: Therefore, the blessed souls demand vengeance of their enemies, although they do not intend it chiefly, because of a zeal for righteousness and an affection of godly love, they resent it, as does God himself at the wickedness of the persecutors, who attack God himself, and seek to hinder his religion, and torment those who worship him; therefore, they would have their malice and power ended. Thus far he. But since Scripture consistently witnesses that the saints in heaven are free from griefs and affections, and live a new life far removed from all pain and disturbance, and have submitted their wills to the will of God, whom they may follow in all things, approving all his judgments, sayings, and doings, and reverencing them: I suppose we need not reason more subtly about this at present, but simply understand that by this figurative speech (as crying is also elsewhere attributed to the blood of martyrs shed) is signified that the blood of the oppressed shall never be forgotten by God, and that before him a just judgment and vengeance is prepared, to be executed in his time against the enemies and contemners of God: but chiefly against the persecutors of the word and the murderers of saints. This is further explained by what follows.

For by the same reasoning that was given in response to the complaints of martyrs, so that we may understand the state and glory of the saints in heaven, who have offered their bodies for God’s testament: and that God has not forgotten the blood spilled, but that he will at length requite those bloodshedders when he sees fit. But where he has reserved this time to himself, when he will reward those who drink the blood, it is not our part to inquire curiously thereof, but rather to be in readiness, that if he will that we also should suffer for the testimony of Jesus Christ, we should run speedily and cheerfully through afflictions unto glory, doubting nothing that we shall be joined to the blessed martyrs in heaven, and that the just judge in that day will render to all the enemies of God, the Church, and God’s word, according to their deeds. And although the time of persecution seems long to the flesh, yet it is here and elsewhere in the scriptures called short. But these things must be seen and considered in parts.

Undoubtedly, the condition of souls in heaven is in every respect most fortunate. This is represented by the white garments. The glory of the blessed, who are now in light and feel nothing of darkness, is signified; of this garment I have spoken before. And it is expressly stated that white garments are given to each of them. For every soul receives its reward; and the body also, at the end of the world, shall receive its own garment. Saint Gregory says that saints as yet have only the blessedness of souls; but at the end of the world, they shall receive two robes or garments, for they shall be clothed with both the perfect joy of souls and the incorruption of bodies. This will be examined more carefully toward the end of Chapter 7, where this passage will be explained more fully. After it was said to the blessed souls that they should rest, therefore they are altogether at peace and feel no disturbance, which will be declared more plainly in Chapter 12. Nevertheless, this may be referred to delay and breathing, as though it were said: It was signified to the souls that they should yet wait and abide. For it follows: yet a little while. Therefore God signifies that after a little time he will deliver his own and punish the adversaries.

And the indication of the time appears to be drawn from the second chapter of Habakkuk, a passage also cited by the Apostle in Hebrews 11. For yet a little while, because he who is to come will come, and will not delay. And the righteous shall live by faith, and so forth. In Isaiah 26, we read these words (after he has shown the resurrection and the final end of the world): “Go, therefore, my people, and enter into your chambers, and shut your doors after you; hide yourselves a little while, even for a moment, until the indignation passes.” And likewise Saint Peter called this entire time of affliction, which is a judgment, a short time, so that we might find comfort in it. 1 Peter 1 and 2 Peter 3.

\index{Judgment, Last}To these matters is joined another, which completes the time until the fulfillment of three companions and brethren who would also be slain for the word of God. Therefore, let us no longer inquire after the end of persecution, or why the Lord delays vengeance, or how long [?]. For we hear that the number of the elect must be accomplished. But since that time is known to God alone, let us not be curious, but let us think of those things that concern our duty, so that if the situation requires it, we may also die bravely for the testament of our Lord God, that we may be associated with our brethren and companions, and have the blessed sight of our redeemer. The number shall be accomplished at the end of the world, at the last judgment. Until then, persecution shall continue; but then assuredly the Lord will requite it, as the prophet Malachi has witnessed in the third and fourth chapters.

Hereof we learn also what we should judge of the holy Martyrs in Heaven, and blessed souls: the same that we learn here of God’s word. Brethren and friends, fellow servants (as also angels, in the 19th and 22nd verses of the Apocalypse) they are expressly called, not Lords and founders. For although the words must be understood of us yet living, there is a relation. For if we be their brethren and fellow servants, they are verily our brethren also, and God’s servants with us, and even fellow servants. Now though we grant that they pray in Heaven, what, I pray you, do they pray for here, but that God would avenge and punish? And what do they obtain? As we read that Christ said to his mother, requiring wine at the marriage, “Woman, what have I to do with thee? My hour is not yet come,” so likewise are the Martyrs here commanded to tarry and abide the time of God appointed. Which we believe that the Saints do. What can we say then of their intercession and praying for sinners unto God? There is one only mediator given, even the Lord Christ; let us go unto him in all our necessities, he alone shall suffice all, and in all.

These things have been spoken up to this point concerning the persecutions of all times, so that in the meantime they have provided the most comforting consolations to all who suffer persecution until the end of the world. They have also silenced curious questions and left us safe and whole in the will of God, in which, resting alone, we may know that it is best for us.

Therefore, it behooves us to gather certain sure sayings, with which to comfort ourselves as with the most certain sentences of God pronounced. First, that God is true and just: and therefore not to neglect his, but to cherish with fatherly care. And if he casts us into any danger or difficulty, that verily shall turn to great profit for the godly. If he shall take us away by torments, that he delivers us from evils, from miseries, and corruption of this world, and renders for the same everlastingness. Secondly, it is certain that God is just and true, that he will requite the wicked after their deserts. Again, if he makes men fortunate in this world, that in deed appertains to their destruction. Where he is slow to punish, that is done through God’s long suffering: but that God recompenses this slowness with the weightiness of the punishment, in case they be incurable. Whereas these things undoubtedly are most certain, what remains but that we should commit ourselves and all ours to the Lord our God? He knows the time and means whereby to avenge his, and to plague his enemies. To him be glory forevermore. Amen.
